\section{Forward expansions}
\label{sec:forward}

Before inverting anything we need the forward function in the normal form $ax^p(1+\sum_ix^{\delta_i}B_i(\log x))$, and the second question \cite{q2} is a forward problem in its own right. This section shows that a large class of expressions has convergent power--log expansions, and establishes rules for propagating a rigorous remainder class through arithmetic and composition.

\subsection{Convergent power--log expansions}

\begin{definition}
\label{def:pl-expansion}
A function $F$ on $(0,w_0)$ has a \emph{convergent power--log expansion} with support in $\beta_0+S$ ($S$ a finitely generated positive semigroup with gaps $\delta_1,\dots,\delta_m$) if there are polynomials $C_\sigma\in\R[L]$, constants $M,K\ge1$, and integers $d_0,d\ge0$ such that
\begin{equation}
 F(w)=\sum_{\sigma\in S}w^{\beta_0+\sigma}C_\sigma(L)\quad\text{for }0<w<w_1,
 \qquad \|C_\sigma\|_1\le M K^{n(\sigma)},\quad
 \deg C_\sigma\le d_0+d\,n(\sigma),
 \label{eq:pl-expansion}
\end{equation}
where $n(\sigma)$ is the least $|\bk|$ with $\bk\cdot\bdelta=\sigma$, $\|P\|_1$ is the sum of the absolute values of the coefficients of $P$, and the series converges absolutely. Thus $|C_\sigma(L)|\le M K^{n(\sigma)}\M(w)^{d_0+dn(\sigma)}$. Zero coefficients satisfy the degree bound by convention. We write $F\in\mathcal E$. The fixed allowance $d_0$ is essential: it admits a polynomial in $\log w$ at the leading exponent, including $\log w$ itself.
\end{definition}

The coefficient bound makes the tail after any cutoff controllable (Lemma~\ref{lem:tail} below), and it is stable under all the operations we need.

\begin{lemma}[Tail of a convergent expansion]
\label{lem:tail}
Let $F\in\mathcal E$ as in \eqref{eq:pl-expansion}, $\delta_*=\min_i\delta_i$, and $H>0$. Then
\begin{equation}
 F(w)-\sum_{\sigma<H}w^{\beta_0+\sigma}C_\sigma(L)=\OO\bigl(w^{\beta_0+H}\M(w)^{d_0+Nd}\bigr),\qquad N=\lceil H/\delta_*\rceil.
 \label{eq:tail}
\end{equation}
More precisely, the omitted part equals $w^{\beta_0+\sigma_*}C_{\sigma_*}(L)+\oo(w^{\beta_0+\sigma_*})$ where $\sigma_*=\min\{\sigma\in S:\sigma\ge H\}$.
\end{lemma}
\begin{proof}
Split the omitted terms into those with $n(\sigma)\le N-1$ and the rest. The first group is finite (Lemma~\ref{lem:finite}) and each term is $\OO(w^{\beta_0+H}\M^{d_0+(N-1)d})$. For the rest, the number of weights whose shortest representation has length $n$ is at most $m^n$, and $\sigma\ge n\delta_*$, so their sum is bounded by
\[
 Mw^{\beta_0}\M^{d_0}\sum_{n\ge N}\bigl(mKw^{\delta_*}\M^d\bigr)^n
 =\OO(w^{\beta_0+N\delta_*}\M^{d_0+Nd}).
\]
Since $N\delta_*\ge H$, this proves the bound. For the refined statement let $\sigma_{**}$ be the next element of $S$ after $\sigma_*$, which exists by local finiteness when $S$ has a generator, and choose $\sigma_*<H'<\sigma_{**}$. Applying the bound at $H'$ shows that everything beyond $\sigma_*$ is $\OO(w^{\beta_0+H'}\M^{d_0+N'd})=\oo(w^{\beta_0+\sigma_*})$. The difference $\sigma_{**}-\sigma_*$ need not be at least $\delta_*$. Finite support with no further terms has zero tail.
\end{proof}

\begin{theorem}[Closure]
\label{thm:closure}
The class $\mathcal E$ contains the constants and $w$ and is closed under sums, products, and the following operations, with the supports and gaps indicated.
\begin{enumerate}[leftmargin=1.8em]
\item If $F\in\mathcal E$ has leading block $cw^\beta$ with a \emph{constant} $c>0$, then $F^a\in\mathcal E$ for every real $a$ and $\log F\in\mathcal E$; their supports are contained in $a\beta+S'$ and $\{0\}\cup S'$ respectively, with block $\log c+\beta L$ at exponent $0$ in the logarithm. Here $S'$ is a finitely generated positive semigroup containing the exponents of $F/(cw^\beta)-1$; the proof establishes that such a semigroup exists.
\item If $F\in\mathcal E$ has no negative-exponent blocks and its exponent-$0$ block is $c+kL$ with $k\in\R$, then $e^F\in\mathcal E$ with leading block $e^cw^k$. The argument $F$ may grow logarithmically.
\item If $F\in\mathcal E$ has $\beta_0\ge0$ and constant exponent-$0$ block $c$, and $\phi$ is analytic at $c$, then $\phi\circ F\in\mathcal E$.
\end{enumerate}
\end{theorem}
\begin{proof}
First note that a positive-valuation expansion $U$ can be represented with support in a positive semigroup based at $0$: if necessary, add its positive leading exponent as a generator and enlarge the constants in \eqref{eq:pl-expansion}. Choose one shortest multi-index $\bk$ for each resulting weight $\sigma$, put $D=d_0+d$, and introduce
\[
 T_{ij}=w^{\delta_i}L^j,\qquad 1\le i\le m,\quad 0\le j\le D.
\]
For $n=|\bk|\ge1$, every $w^\sigma L^s$ with $s\le d_0+dn\le Dn$ is a product of $n$ such parameters: distribute the $s$ logarithmic factors among the $n$ generator factors, each receiving at most $D$. Fix one such distribution for every monomial. The coefficient-norm bound makes the resulting series analytic in a sufficiently small polydisc, since its terms of degree $n$ have total absolute coefficient at most $M(mK)^n$. Analytic composition with $(1+U)^a$, $\log(1+U)$, $e^U$, or $\phi(c+U)$ is therefore legitimate. Cauchy estimates give exponential coefficient bounds in the new parameters, and their logarithmic degree grows at most linearly in parameter degree.

Regrouping by weight preserves bounds of the form \eqref{eq:pl-expansion}: a representation of weight $\sigma$ has length at most $\sigma/\delta_*$, whereas its shortest length is at least $\sigma/\delta_{\max}$; their ratio is bounded, and the number of parameter monomials of a given total degree is at most exponential. Enlarge $M,K,d$ to absorb both factors.

For sums and products, use the union of the generator sets and, for a sum of differently shifted supports, the positive difference between their shifts as an additional generator. Coefficient norms are subadditive and submultiplicative, so the same bounds apply. If leading blocks cancel, rebasing at the next nonzero exponent still admits finitely many positive generators: the multi-indices above that exponent form an upward-closed subset of $\N^m$, with finitely many coordinatewise minimal elements. Subtracting the new leading exponent from the weights of those elements supplies finitely many nonnegative shifts, together with the original generators. This also justifies normalizing $F$ by its actual leading monomial in (1). Separate the finite leading polynomial before applying the small-parameter construction; in (2) separate $c+kL$ to obtain the factor $e^cw^k$.
\end{proof}

The theorem supplies sufficient closure conditions, and every truncation is a genuine asymptotic expansion by Lemma~\ref{lem:tail}. They are not necessary conditions for every individual expression: for example $\sqrt{L^2}=-L$ for $L<0$, and polynomials may be applied to an unbounded logarithm. In general, non-integer powers of a nonconstant leading polynomial, its logarithm, exponentials of negative-power arguments, and oscillatory functions of an unbounded logarithm require larger scales (Section~\ref{sec:boundaries}).

\begin{corollary}[Laurent and Puiseux composition]
If $U\in\mathcal E$ has a positive constant leading coefficient and positive valuation, and a selected real branch of $\phi(c+v)$ has a convergent Puiseux expansion $v^{n_0/d}A(v^{1/d})$ for $v\to0^+$, with $d\ge1$ an integer and $A$ analytic at $0$, then $\phi(c+U)\in\mathcal E$. The same holds from below using $\phi(c-v)$.
\end{corollary}
\begin{proof}
Apply the closure theorem to $U^{1/d}$, the analytic function $A$, and $U^{n_0/d}$. Negative $n_0$ includes poles. The side of approach fixes the real branch of a ramified function.
\end{proof}

Variable powers also reduce to the closure theorem through $F^G=\exp(G\log F)$ when $F>0$ and the resulting exponential argument satisfies its hypotheses; this includes $w^w=\exp(w\log w)$. For a nonzero leading block $w^\beta P(L)$, the eventual sign is the sign of $(-1)^{\deg P}[L^{\deg P}]P$, so absolute values reduce to multiplication by that sign whenever it is provable.

\subsection{Modulus of a complex-valued expansion}
\label{sec:modulus}

The sign reduction just stated is a statement about \emph{real} coefficient
polynomials. Write $\mathcal E_\C$ for the class of Definition~\ref{def:pl-expansion}
with $C_\sigma\in\C[L]$ instead of $\R[L]$. It is closed under sums and
products by the same argument, since \eqref{eq:pl-expansion} bounds
$\|C_\sigma\|_1$ and never uses reality. For $F\in\mathcal E$ with a nonzero
leading block the eventual sign gives $|F|=\pm F$. For
$F\in\mathcal E_\C\setminus\mathcal E$ that rule is false, and its failure is
invisible in the leading block.

\begin{example}[A positive leading coefficient does not license the sign rule]
\label{ex:modulus-cancellation}
Let $a\in\C$ satisfy $a^2=-1$, and let $w\to0^+$ be real. Both $1+aw$ and
$1-aw$ have leading block $1$, with positive constant leading coefficient.
Replacing each modulus by its argument and adding gives the identically zero
expression $(1+aw)+(1-aw)-2$. The true value is
\[
 |1+aw|+|1-aw|-2=2\sqrt{1+w^2}-2=w^2-\tfrac14w^4+\OO(w^6),
\]
whose leading term is $w^2$. The whole expansion is lost, and the loss is
concealed: the two imaginary parts cancel, so the returned coefficients are
real. Reality of the \emph{result} of a compound calculation is therefore no
evidence that a modulus taken inside it was evaluated correctly.
\end{example}

The correct construction stays on the real coordinate. Because $w$ and
$L=\log w$ are real, conjugation acts on $\mathcal E_\C$ coefficientwise:
$\overline F=\sum_\sigma w^{\beta_0+\sigma}\overline{C_\sigma}(L)$, where
$\overline{C_\sigma}$ conjugates the coefficients of the polynomial.

\begin{proposition}[Modulus on a real coordinate]
\label{prop:modulus}
Let $F\in\mathcal E_\C$ and put $Q=F\overline F$. Then $Q\in\mathcal E$, $Q\ge0$
pointwise, and its blocks are the ordered convolution
$Q_\sigma=\sum_{\tau+\tau'=\sigma}C_\tau\overline{C_{\tau'}}$ at exponents
$2\beta_0+\sigma$. If the leading block of $F$ is a nonzero constant $c\in\C$,
then $Q$ has leading block $|c|^2w^{2\beta_0}$ and
\[
 |F|=Q^{1/2}\in\mathcal E,
\]
with leading block $|c|w^{\beta_0}$ and support contained in $\beta_0+S'$ for a
finitely generated positive semigroup $S'$.
\end{proposition}
\begin{proof}
The convolution is the product rule of Theorem~\ref{thm:closure}, which holds
verbatim over $\C[L]$ as noted above, applied to $F$ and $\overline F$; they
have the same support and coefficient bounds. Reindexing by
$(\tau,\tau')\mapsto(\tau',\tau)$ and conjugating fixes each $Q_\sigma$, so
$Q_\sigma\in\R[L]$ and $Q\in\mathcal E$, this time by the theorem itself.
Pointwise $Q(w)=|F(w)|^2\ge0$. Since $\tau+\tau'=0$ with $\tau,\tau'\ge0$
forces $\tau=\tau'=0$, no cancellation can occur in the leading block, and a
nonzero constant leading coefficient gives $c\overline c\,w^{2\beta_0}=|c|^2w^{2\beta_0}$
with $|c|^2>0$. Part~(1) of Theorem~\ref{thm:closure} then applies to $Q$ with
$a=1/2$ and returns the positive branch, so $|F|$ and $Q^{1/2}$ are nonnegative
with the same square, hence equal.
\end{proof}

The constant leading coefficient is not a technicality. If the leading block is
a nonconstant polynomial in $L$, the square root of $|C_{0}|^2$ need not be
polynomial in $L$, and the modulus leaves the scale: for $F=L+a$ with $a^2=-1$
one has $|F|=\sqrt{L^2+1}$, which is not a power--log expansion in the sense of
Definition~\ref{def:pl-expansion} and belongs to the larger scales of
Section~\ref{sec:boundaries}. The earlier observation that $\sqrt{L^2}=-L$ for
$L<0$ shows that a provable sign can still rescue such a case; here there is no
sign to prove.

\begin{remark}[Error transport and cost]
\label{rem:modulus-transport}
The reverse triangle inequality $\bigl||F|-|T|\bigr|\le|F-T|$ transports an
absolute remainder through the modulus unchanged: if $F-T=\OO(w^P\M^D)$ then
$|F|-|T|=\OO(w^P\M^D)$. This bounds magnitude only. It identifies no leading
term, fixes no sign, and supplies no derivative information; Section~\ref{sec:series-calculus}
shows that a modulus does not transport a classical derivative contract even
when it transports the magnitude faithfully. A retained support of $r$ blocks
admits $r^2$ ordered coefficient products in $Q$. For $\tau\ne\tau'$ the two
products $C_\tau\overline{C_{\tau'}}$ and $C_{\tau'}\overline{C_\tau}$ are
conjugates, so their contribution is $2\operatorname{Re}\bigl(C_\tau\overline{C_{\tau'}}\bigr)$,
while the diagonal product $C_\tau\overline{C_\tau}=|C_\tau|^2$ is already real
and enters once. Only the $r(r+1)/2$ unordered pairs therefore need to be
formed. For $r=64$ that is $2080$ products instead of $4096$.
\end{remark}

\subsection{Propagation of finite-order errors}
\label{sec:evaluator}

A finite-order expansion of a function $F$ consists of a jet $T$ and a remainder assertion
\begin{equation}
 F=T+\OO\bigl(w^P\M(w)^D\bigr),\qquad P\in\R\cup\{+\infty\},
 \label{eq:triple}
\end{equation}
where $P=+\infty$ denotes an exact finite expression. The following rules construct expansions of compound functions from those of their constituents.
\begin{itemize}[leftmargin=1.6em]
\item For a sum, add the jets and use the coarser of the two remainder classes. If a block is discarded at the boundary exponent, include its whole logarithmic degree in that remainder.
\item For a product $(T_1+R_1)(T_2+R_2)$, the unknown part is $R_1T_2+R_2T_1+R_1R_2$. Apply \eqref{eq:O-arithmetic} to all three contributions, and to every explicit product block removed by truncation. A negative leading exponent in either factor can reduce the absolute precision.
\item Suppose $F=cw^\beta(1+U)$ with $c>0$ and $U\to0$. An absolute error $\OO(w^P\M^D)$ in $F$ is a relative error $\OO(w^{P-\beta}\M^D)$ in $U$. For fixed real $a$, expand $(1+U)^a$ and multiply by $c^aw^{a\beta}$; the propagated input error is $\OO(w^{P+(a-1)\beta}\M^D)$. Integer powers also follow by finite multiplication, without a positivity assumption.
\item With the same positive leading monomial, $\log F=\log c+\beta L+\log(1+U)$. The propagated error is measured at relative exponent $P-\beta$.
\item If $F=c+kL+U$, where $k\in\R$, $U\to0$ and the input uncertainty tends to zero, then $e^F=e^cw^ke^U$. The input error is multiplied by the exact factor $e^cw^k$.
\item For a function $\phi$ analytic at $c$ and a vanishing increment $U$, use $\phi(c+U)=\sum_{j=0}^{n}\phi^{(j)}(c)U^j/j!+\OO(U^{n+1})$. Laurent and Puiseux expansions on a fixed side reduce to powers of the signed increment and analytic composition.
\end{itemize}
These analytic unit substitutions require their propagated relative uncertainties to tend to zero. All unit-series truncations are finite by positive valuation. The error from their explicit truncation is combined with the propagated uncertainty of their arguments. If cancellation conceals the leading increment needed for a pole or a ramified branch, more terms of that increment must first be determined.

\begin{lemma}[Nonlinear logarithmic bound at the precision boundary]
\label{lem:unit-boundary-degree}
Let $U=\sum_{\beta>0}w^\beta P_\beta(L)$ be a nonzero finite jet,
$\nu=\val U>0$, and $d=\max_\beta\deg P_\beta$. Suppose $F$ is
real-valued and $F=U+\OO(w^P\M^D)$ with $P>0$ and $D\in\N$. An exact argument $F=U$
is also allowed, with $P=+\infty$. Let $\Phi$ be real analytic near
$0$, and let $H>0$ be finite. Put
\[
 c=\min(P,H),\qquad N=\lceil c/\nu\rceil,\qquad
 J=\left[\sum_{j=0}^{N-1}\frac{\Phi^{(j)}(0)}{j!}U^j\right]_{<c},
\]
where the brackets retain complete blocks of exponent less than $c$.
Then
\begin{equation}
 \Phi(F)=J+\OO(w^c\M^K),\qquad
 K=\begin{cases}
 Nd,&H<P,\\
 \max(D,Nd),&P\le H.
 \end{cases}
 \label{eq:unit-boundary-degree}
\end{equation}
\end{lemma}
\begin{proof}
The finite jet satisfies $U=\OO(w^\nu\M^d)$ and tends to zero.
Taylor's theorem gives a tail $\OO(U^N)$, which is
$\OO(w^{N\nu}\M^{Nd})\subseteq\OO(w^c\M^{Nd})$.
Every block discarded from the finite Taylor polynomial has exponent
at least $c$ and logarithmic degree at most $(N-1)d$, so it obeys the
same bound. Both $F$ and $U$ tend to zero, and $\Phi'$ is bounded on
a neighborhood containing them. The mean value theorem therefore gives
$\Phi(F)-\Phi(U)=\OO(w^P\M^D)$. If $H<P$, Lemma~\ref{lem:dominance}
absorbs this into $\OO(w^H)$; if $P\le H$, its degree must be combined
with the nonlinear truncation degree at the same power. For $F=U$ this
last contribution is absent.
\end{proof}

In particular, retaining more powers than the input precision permits
does not remove the logarithms generated by nonlinear truncation. For
$F=wL+\OO(w^3)$ and $H\ge3$,
\[
 \log(1+F)=wL-\tfrac12w^2L^2+\OO(w^3\M^3).
\]
The cubic Taylor term has degree three although the inherited degree is
zero. The bound uses the largest degree of \emph{all} retained blocks:
$U=w+w^2L^7$, $P=3$, $D=0$ gives the valid conservative degree $21$,
without asserting it is minimal. If $U=0$, boundedness of $\Phi'$ instead
directly gives $\Phi(F)=\Phi(0)+\OO(w^P\M^D)$; for $F=0$ this is exact.

\begin{proposition}[Soundness of compositional error propagation]
\label{prop:soundness}
Suppose the initial expansions satisfy their stated branch and remainder hypotheses. Every finite composition of the rules above satisfies its resulting assertion \eqref{eq:triple}.
\end{proposition}
\begin{proof}
Induct on the number of operations. For sums and products, use \eqref{eq:O-arithmetic} and include the explicit discarded blocks. Lemma~\ref{lem:unit-boundary-degree} combines the Taylor truncation error with the argument error of an analytic unit function; Lemma~\ref{lem:tail} supplies finite-order input bounds for convergent expansions. Factoring the leading monomial gives the stated shifts for powers and logarithms; the same factorization reduces a Laurent or Puiseux composition to the preceding rules.
\end{proof}

Each initial analytic remainder estimate is a hypothesis of this proposition. Knowing the degree of one omitted coefficient does not by itself bound an arbitrary unknown tail. For a convergent expansion in the class $\mathcal E$, Lemma~\ref{lem:tail} supplies the needed complete-tail assertion. For other functions, an independent finite-order asymptotic theorem is required.

\subsection{The second question}

For $F(x)=(1+x+x^{\sqrt2})^{\sqrt2}$ at $x\to+\infty$ the local variable is $w=1/x$ and
\[
 F=w^{-2}\bigl(1+w^{\sqrt2-1}+w^{\sqrt2}\bigr)^{\sqrt2}=w^{-2}\sum_{k\ge0}\binom{\sqrt2}{k}\bigl(w^{\sqrt2-1}+w^{\sqrt2}\bigr)^k ,
\]
an instance of item (1) of Theorem~\ref{thm:closure} with gaps $\sqrt2-1$ and $\sqrt2$. The exponents of $x$ are $2-(\sqrt2-1)k_1-\sqrt2k_2$; ordering the semigroup elements gives \eqref{eq:answer3}, and the first omitted exponent is $7-5\sqrt2$ (multi-index $(5,0)$). The tail estimate therefore gives the remainder $\OO(x^{7-5\sqrt2})$ in \eqref{eq:answer3}. Every larger finite cutoff is treated by the same multinomial expansion and ordering of the positive-gap semigroup.
